\section{原作者完整算法流程}
\label{sec:03_04}

Großmann等提出的完整求解路径由精确归约、残余问题求解和逆序展开组成\cite{GrossmannJRS2026}。归约先识别可以固定、删除或暂缓决策的任务；残余问题上的贪心与局部搜索负责寻找可行选择；逆序展开再把这些选择恢复为原始任务编号。本节说明这一公开算法的数学依据及实际调用配置，不将其作为本书提出的新算法。原论文允许在归约后调用整数线性规划求解器，也允许采用启发式方法。已有完整流程实验采用归约、$S,\mathrm{pin}$静态贪心、ILS15和逆序展开，其中没有SCIP求解残余问题的步骤。算法设计中的规则、冻结源码能够检测的结构以及历史输出的有效性，分别具有不同的证据含义。

\subsection{归约状态、残余核与目标偏置}

原始输入记为$H_0=(V_0,E_0)$。资源集合$V_0$中的每个元素表示一种受容量限制的资源，任务索引集合$E_0$中的每个元素表示一个候选任务。令$m=|V_0|$、$n=|E_0|$；任务$e\in E_0$使用非空资源集合$R_e\subseteq V_0$，具有无量纲正收益$w_e$。实际适配器采用整数权重，因此输入满足$w_e\in\mathbb{Z}_{>0}$。本节的详细归约条件限定于已有组合实验使用的单位容量，即$b_v=1$，$v\in V_0$。容量、任务数和资源数均为无量纲计数。

对任意$M\subseteq E_0$，其收益为$W_0(M)=\sum_{e\in M}w_e$。原始问题为
\begin{equation}
 z_0^*=\max_{M\subseteq E_0}
 \left\{W_0(M):
 \sum_{e\in M}\mathbb{I}_{\{v\in R_e\}}\leq1,
 \ \forall v\in V_0\right\},
 \label{s0304:original_problem}
\end{equation}
其中$\mathbb{I}_{\{P\}}$在命题$P$成立时取$1$，否则取$0$。一个任务必须同时获得其全部关联资源，因而两项任务只要共享一个资源，就不能同时选择。

归约过程中需要区分四个对象。第一，固定集合$F$记录已由安全归约确定选入的当前任务；第二，残余核$K=(V_K,E_K)$保存尚需决策的资源与任务；第三，历史栈$\mathcal{H}$按发生顺序保存折叠、合并及恢复所需的对应关系；第四，目标偏置$C\geq0$累计已经从残余目标中移出的收益。这里的核指经过有限轮归约后留下的问题，不保证所有归约规则均已不能继续应用。

残余任务$a\in E_K$具有当前资源集合$R_a^K\subseteq V_K$和当前权重$\omega_a>0$。下标$a$可以指原始任务，也可以指折叠产生的临时任务，因而$E_K$不能直接当作$E_0$的子集。固定集合也可能包含临时任务，必须结合历史栈才能解释其原始含义。$C$、$\omega_a$和原始收益使用相同的无量纲尺度。初始时$F=\varnothing$、$K=H_0$、$\mathcal{H}$为空且$C=0$。这些对象构成解释归约的数学记账方式，不要求源码内部逐一设置同名字段。

在这一表示中，固定选入的任务从残余决策集合中移除，其占满资源上的其他任务也被排除；尚可使用的资源保持单位剩余容量。实现可以通过匹配状态、阻挡标记和剩余容量保存相同信息，而不必立即物理删除全部对象。以下数学规则中的关联集合均指当前仍需决策的任务，不能把实现中已经阻挡但尚存于容器的任务重新作为可选项。

定义残余收益$W_K(S)=\sum_{a\in S}\omega_a$，其中$S\subseteq E_K$满足当前单位容量约束。若每次归约均满足精确性条件，并完整保存恢复关系，则有
\begin{equation}
 \begin{aligned}
 z_0^*&=C+z_K^*,\\
 z_K^*&=\max_{S\subseteq E_K}
 \left\{W_K(S):
 \sum_{a\in S}\mathbb{I}_{\{v\in R_a^K\}}\leq1,
 \ \forall v\in V_K\right\}.
 \end{aligned}
 \label{s0304:optimum_offset}
\end{equation}
星号表示相应问题的全局最优值。偏置只累计已经移出残余目标的收益，不能在恢复时再重复加入固定任务的收益。

设$\mathcal{U}_{\mathcal{H},F}$表示按照历史记录恢复固定任务和折叠任务的操作。对一个可行残余解$S$，安全归约及正确恢复给出
\begin{equation}
 M=\mathcal{U}_{\mathcal{H},F}(S)\subseteq E_0,
 \qquad W_0(M)=C+W_K(S).
 \label{s0304:solution_offset}
\end{equation}
式\eqref{s0304:optimum_offset}是精确归约的最优值关系；式\eqref{s0304:solution_offset}则说明某个实际解如何记账。如果$S$由启发式得到，只能由可行性得出$W_K(S)\leq z_K^*$，不能因为归约具有精确性，就把恢复后的$M$称为全局最优解。这些关系的应用还以实际执行的变换满足相应条件为前提，不能由函数名称中的归约二字直接推定。

\subsection{删除资源、固定任务与支配删除}

为说明各条规则，在当前残余核上定义资源关联任务集、开邻域和闭邻域：
\begin{align}
 E_K(v)&=\{a\in E_K:v\in R_a^K\}, &&v\in V_K,
 \label{s0304:incident}\\
 N_K(a)&=\{f\in E_K\setminus\{a\}:R_f^K\cap R_a^K\ne\varnothing\},
 &&a\in E_K,
 \label{s0304:open_neighborhood}\\
 N_K[a]&=N_K(a)\cup\{a\}, &&a\in E_K.
 \label{s0304:closed_neighborhood}
\end{align}
邻域表示资源竞争关系，不表示执行顺序。以下条件均为足以保证相应操作安全的条件，不是必要条件。

\paragraph{充裕资源删除AV。}
若资源$v$的关联任务数不超过其剩余容量，则无论这些任务中选取哪些，该资源都不会成为限制。在上述已经消去固定占用的残余表示中，条件为$|E_K(v)|\leq1$。因此可以从当前资源集合及所有关联任务的资源列表中删除$v$。实际代码比较的是剩余容量与当前可匹配关联任务数，不能将原始度数与更新后的剩余容量混用。

若某个任务$a$在删除充裕资源后成为空资源任务，则它与全部剩余任务兼容。由于$\omega_a>0$，可将$a$移入固定集合，从残余任务集中移除，并更新$C\leftarrow C+\omega_a$。原始资源关联仍须保存在恢复和核验所用的输入中；从归约问题中删除资源，并不表示原始任务不再使用该资源。

\paragraph{邻域删除NR。}
对每个$v\in R_a^K$，定义通过该资源与$a$竞争的最大单项收益
\begin{equation}
 q_v(a)=\max\bigl(\{\omega_f:f\in E_K(v)\setminus\{a\}\}\cup\{0\}\bigr).
 \label{s0304:nr_resource_bound}
\end{equation}
$q_v(a)\geq0$是无量纲收益界，加入$0$用于处理没有其他关联任务的情况。单位容量下，以下条件足以固定$a$：
\begin{equation}
 \omega_a\geq\sum_{v\in R_a^K}q_v(a).
 \label{s0304:nr_condition}
\end{equation}
任何可行解在每个$v\in R_a^K$上至多选择一个任务。若该解没有选择$a$，删除其中与$a$冲突的任务后，损失的总收益不超过式\eqref{s0304:nr_condition}右端，再加入$a$便不会降低收益。同一邻接任务可能经多个资源被重复计入右端，这只会使判定更保守，不破坏其充分性。

满足条件后，将$a$固定选入，删除$a$及$N_K(a)$，并令$C\leftarrow C+\omega_a$。这里的固定含义是至少存在一个最优解包含$a$，并不表示所有最优解都必须包含它。原论文的一般容量规则使用每个资源处排除$a$后的第$b_v$大邻边权重；在残余状态中相应使用剩余容量。该规则不是前$b_v$项权重之和。本节仅展开单位剩余容量时的最大权重形式。

\paragraph{加权孤立边删除WIER。}
若$a$在其闭邻域中最重，而且闭邻域内任意两个不同任务均冲突，即
\begin{align}
 \omega_a&\geq\omega_f,
 &&\forall f\in N_K(a),
 \label{s0304:wier_weight}\\
 R_f^K\cap R_g^K&\ne\varnothing,
 &&\forall f,g\in N_K[a],\ f\ne g,
 \label{s0304:wier_clique}
\end{align}
则可以固定$a$并删除其闭邻域，偏置增加$\omega_a$。由于全部资源容量为$1$，式\eqref{s0304:wier_clique}保证可行解在该闭邻域中至多选择一项任务。若选择了邻居$f$，可用收益不低的$a$替换；闭邻域之外的任务均不与$a$冲突。如果闭邻域中一项也未选，则可以直接追加$a$。

式\eqref{s0304:wier_clique}要求的是每对任务各自存在共享的单位容量资源。这些共享资源可以因任务对而不同，并不要求闭邻域内所有任务共同使用同一个资源。后一条件更强，会遗漏满足原规则充分条件的可归约结构。

\paragraph{加权支配删除WD。}
对不同任务$a,f\in E_K$，若
\begin{equation}
 R_a^K\subseteq R_f^K,
 \qquad \omega_a\geq\omega_f,
 \qquad R_a^K\ne\varnothing,
 \label{s0304:wd_condition}
\end{equation}
则可以删除$f$，偏置不变。非空条件在单位容量下保证$a$与$f$至少共享一个容量为$1$的资源，二者不能同时选入。任何包含$f$的可行解都可以用$a$替换$f$：替换后的资源需求是原需求的子集，收益也不会下降。资源集合相同且权重相同的两项任务允许保留其中一个代表，但必须保留原始编号对应关系。

实际WD代码仅把权重严格低于$a$的任务作为待删除候选，并进一步限制候选大小和扫描范围。因而它检测的是上述安全条件的一部分，不包含理论上允许的全部等权支配情形。某个可支配任务未被删除，不等于支配性质不成立，也不能据此认为残余核已经穷尽全部安全删除。

\subsection{暂缓选择的边折叠与孪生合并}

固定规则直接确定某项任务进入一个最优解，WD则排除一项可被替代的任务。折叠处理的是另一种情况：局部结构已能缩小选择范围，但具体应选哪一组仍取决于外部任务。

\paragraph{加权边折叠WEF的安全条件。}
取当前任务$a$，令$N=N_K(a)$为其完整开邻域，且$N\ne\varnothing$。要求邻域中的任务彼此独立，即
\begin{equation}
 R_f^K\cap R_g^K=\varnothing,
 \qquad \forall f,g\in N,\ f\ne g.
 \label{s0304:wef_independence}
\end{equation}
每个$f\in N$都与$a$共享单位容量资源，因此选择$a$会排除全部$N$；式\eqref{s0304:wef_independence}则允许$N$中的全部任务同时选择。记
\begin{equation}
 W_N=\sum_{f\in N}\omega_f,
 \qquad \omega_{\min}=\min_{f\in N}\omega_f.
 \label{s0304:wef_neighbor_weight}
\end{equation}
若进一步满足
\begin{equation}
 W_N>\omega_a>W_N-\omega_{\min},
 \label{s0304:wef_weight}
\end{equation}
则可以把$a$与$N$折叠为一个新任务$h$：
\begin{align}
 R_h&=\bigcup_{f\in N}R_f^K,
 &\omega_h&=W_N-\omega_a>0,
 \label{s0304:wef_new_edge}\\
 E_{K'}&=\bigl(E_K\setminus(N\cup\{a\})\bigr)\cup\{h\},
 &C'&=C+\omega_a.
 \label{s0304:wef_update}
\end{align}
这里$K'$为本次折叠后的残余问题，$C'$为更新后的偏置。历史栈保存$a$、$N$、$h$及其对应关系，以便后续恢复。

条件\eqref{s0304:wef_weight}为何有效，可以从任何严格子集$Q\subsetneq N$的收益看出：
\begin{equation}
 \sum_{f\in Q}\omega_f\leq W_N-\omega_{\min}<\omega_a.
 \label{s0304:wef_proper_subset}
\end{equation}
若一个最优解选择了非空严格子集$Q$，将$Q$替换成$a$会增加收益；这一替换不会与邻域之外的任务冲突，因为$N$已经是$a$的完整邻域。若既未选择$a$也未选择任何邻居，则可直接加入$a$。所以可以把最优决策限制为选择$a$或选择整个$N$，而无需保留只选部分邻居的状态。

新任务$h$的资源并集保留了选择整个$N$对外部任务产生的全部冲突。其权重记录选择$N$相对于选择$a$的额外收益。对折叠后可行解$S'$，一次逆变换为
\begin{equation}
 S=\begin{cases}
 (S'\setminus\{h\})\cup N, &h\in S',\\
 S'\cup\{a\}, &h\notin S'.
 \end{cases}
 \label{s0304:wef_unfold}
\end{equation}
选中$h$时恢复全部邻居；未选中$h$时恢复中心任务$a$。两种情况的恢复收益均等于折叠后收益加$\omega_a$。因此折叠权重既不是原始任务收益，也不是已经确定获得的收益，不能直接用于报告原始问题目标值。

\paragraph{实际WEF的检测范围与等权边界。}
冻结源码只检查当前资源数恰为$2$的中心任务。有效折叠还要求中心相关资源的剩余容量为$1$、可匹配关联任务数为$2$，最终收集到恰好两个不同且独立的邻居$f,g$。代码采用的权重条件为
\begin{equation}
 \omega_f\leq\omega_a,\qquad
 \omega_g\leq\omega_a,\qquad
 \omega_f+\omega_g>\omega_a.
 \label{s0304:wef_implementation_weight}
\end{equation}
这些检查在前置AV清理后的当前结构上进行，不能直接用原始任务的资源数或原始资源度数替代。

当$N=\{f,g\}$时，理论条件\eqref{s0304:wef_weight}中的严格下界为$\omega_a>\max\{\omega_f,\omega_g\}$，而式\eqref{s0304:wef_implementation_weight}允许中心与较重邻居等权。这一区别不应写成完全相同的判据。在该受限结构下，若某个最优解仅选择一个邻居，以$a$替换该邻居仍然可行，收益不下降。因此即使发生等权，也至少存在一个只在$a$和完整邻域$N$之间选择的标准形最优解，仍可采用式\eqref{s0304:wef_new_edge}至式\eqref{s0304:wef_unfold}的变换。等权时不能声称替换严格增加收益，也不能声称所有最优解原本都已具有这一标准形。

上述论证同时说明了条件的适用范围。两个邻居时，逐个邻居不重于中心即可控制所有非空严格子集；邻居数增加后，逐项比较不再控制多个邻居组成的严格子集，不能将式\eqref{s0304:wef_implementation_weight}直接推广为任意邻域的折叠条件。

\paragraph{加权孪生WT的数学变换。}
设$a_1,a_2\in E_K$为两个不同且不相邻的任务，满足
\begin{equation}
 R_{a_1}^K\cap R_{a_2}^K=\varnothing,
 \qquad N_K(a_1)=N_K(a_2)=N.
 \label{s0304:wt_structure}
\end{equation}
要求共同邻域$N$独立。若$N$非空，再定义合并权重
\begin{equation}
 \Omega=\omega_{a_1}+\omega_{a_2},
 \qquad \Omega>W_N-\omega_{\min}.
 \label{s0304:wt_weight}
\end{equation}
$\Omega$为同时选择两个孪生任务的无量纲总收益。由于二者互不冲突，且冲突邻居完全相同，选择其中一个而未选择另一个时，总可以把另一个追加进去。因此可以把二者作为一个整体处理。

在单位容量下，可用临时任务$t$代表这一整体，令$\omega_t=\Omega$，并取$R_t=R_{a_1}^K$。只保留一个孪生任务的资源集合仍能保存全部对外冲突：任何与$a_2$冲突的外部任务也在$a_1$的邻域中。取二者资源并集也具有相同的对外冲突关系，这里采用论文正文的代表集合写法。历史栈必须记录$t$被选中时恢复$a_1$和$a_2$，而不是只恢复其中一个。

合并后分两种情况处理。当$\Omega\geq W_N$时，删除任何已选邻居的总收益损失都不超过$W_N$，故可按邻域交换论证固定$t$，恢复时选入两个孪生任务。这属于NR的收益支配逻辑，但使用的是完整邻域总收益界，不要求它等于式\eqref{s0304:nr_condition}中可能重复计数的逐资源上界。当$W_N>\Omega>W_N-\omega_{\min}$时，对$t$和$N$应用WEF：新折叠任务的资源集合为$\bigcup_{f\in N}R_f^K$，权重为$W_N-\Omega$，偏置增加$\Omega$。恢复时，选中折叠任务对应选择$N$，未选中对应选择$a_1,a_2$。如果$N$为空，则两个正权重任务均可直接固定，无需计算空集合的最小权重。

以上WT步骤依据论文正文的合并及交换含义展开，描述的是满足条件时安全的数学变换。它依赖完整邻域相同、孪生任务互不相邻和单位容量等条件，不能仅凭两项任务有部分共同邻居便执行合并；它也不意味着当前源码已经成功识别并执行了这些变换。

\subsection{实际归约顺序、默认参数与实现覆盖范围}

适配器首先按原始输入建立任务资源列表，必要时对列表排序，设置资源容量和随机种子。随后调用公开代码的非配置式归约驱动器。驱动器先按权重组织关联信息，每轮依次调用
\begin{equation}
 \mathrm{AV}\ \to\ \mathrm{NR}\ \to\ \mathrm{WIER}\ \to\ \mathrm{WD}
 \ \to\ \mathrm{AV}\ \to\ \mathrm{WEF}\ \to\ \mathrm{AV}\ \to\ \mathrm{WT}.
 \label{s0304:reduction_order}
\end{equation}
AV在一轮中出现多次，是因为前面的删除与折叠可能使某些资源不再构成限制。驱动器依据节点变化时间戳判断是否继续下一轮，并将轮数上限设为$10$。因此实际执行至多$10$轮，可能提前停止；达到上限也不能解释为已经获得不可继续归约的核。式\eqref{s0304:reduction_order}表示函数调用次序，不保证每个被调用函数均成功应用了一条归约。

原论文表3列出了控制检测开销的默认阈值。适配器没有逐项传入这些规则参数，实际采用冻结源码函数的默认值。二者的对应关系见表\ref{s0304:defaults}。表中大小、数量和轮数均为无量纲整数，默认参数与循环实际处理次数需要分别理解。

\begin{table}[htbp]
 \centering
 \caption{归约检测参数：论文表3与冻结源码默认值}
 \label{s0304:defaults}
 \begin{tabularx}{\linewidth}{llccX}
  \toprule
  规则 & 参数含义 & 论文表3 & 源码 & 解释 \\
  \midrule
  NR & 中心边大小 & $10$ & $10$ & 超过该大小的候选跳过。 \\
  WIER & 中心边大小 & $8$ & $8$ & 限制中心资源关联项数。 \\
  WIER & 团大小 & $80$ & $80$ & 限制收集的闭邻域大小。 \\
  WD & 子边大小 & $6$ & $6$ & 限制拟保留子边的资源数。 \\
  WD & 候选参数 & $6$ & $6$ & 候选加入后再检查是否超过参数。 \\
  WD & 检查参数 & --- & $3$ & 使用后增计数与严格大于比较。 \\
  WT & 中心边大小 & $4$ & $4$ & 用于孪生候选的初步筛选。 \\
  全部 & 归约轮数上限 & $10$ & $10$ & 适配器也显式传入该上限。 \\
  \bottomrule
 \end{tabularx}
\end{table}

例如，WD的候选参数为$6$，但代码先加入一个候选，再判断候选数是否大于$6$，所以列表可以收集到$7$项。其检查参数为$3$，但计数器从$0$开始，采用后增计数并判断旧值是否大于$3$，故可有$4$个通过此前筛选的子边进入后续检查。这些都是具体比较顺序的结果，不能将参数值直接称为精确检查次数。WD还仅收集严格较轻的候选，并经过资源包含关系复核；参数限制影响检测范围，不替代支配条件本身。

冻结源码中的WT存在可由控制流直接识别的逻辑缺陷。首先，候选排序比较器包含对同一个整数键的严格自比较，该比较恒为假，因而整个比较器不能实现预期的按邻域大小和键值分组排序。其次，对相同长度的非空邻域逐项判断是否相等时，另一邻域的迭代器在比较前先前移一步；到最后一项时，该迭代器已经到达末尾，导致相同邻域标志变为假。因此，这条非空等长邻域识别路径会拒绝候选，不能到达后续的孪生固定或折叠分支。调用WT函数与实际成功执行WT归约，在该冻结版本中必须明确区分。

这一结论来自冻结源码的静态控制流，不是性能消融结果。已有公开源码和历史实验均保留原状，未通过修复WT并重跑来测量其影响，因而不能把这一缺陷归因为后述$5$次失败，也不能据此解释具体收益差或耗时差。该实现路径的运行记录不能作为所有理论归约规则均已得到完整实现验证的证据。未执行某条可用归约通常会留下更多残余决策，但当前材料不足以量化这一点对启发式搜索的实际影响。

\subsection{静态贪心与ILS的连接}

归约结束后，适配器对已经确定的匹配部分设置精确标记，再进行$S,\mathrm{pin}$静态贪心。对当前非空残余任务$a$，其排序分数为
\begin{equation}
 s(a)=\frac{\omega_a}{|R_a^K|},
 \qquad a\in E_K,\quad |R_a^K|\geq1.
 \label{s0304:static_pin_score}
\end{equation}
分子是当前任务权重，分母是当前资源关联项数，故$s(a)$为无量纲排序值。它使用归约后的资源集合与权重；临时折叠任务的分子可以是收益差额。排序在贪心阶段保持静态，但每次接收任务仍须检查当前资源容量。空资源任务由归约阶段固定处理，式\eqref{s0304:static_pin_score}不用于零分母情形。

这一分数不同于第\ref{sec:03_03}节诊断和压力实验所用的度乘积优先级。即使两种流程都使用Greedy和ILS，初始化的选择次序、输入结构及可撤销任务集合也可能不同，不能仅凭共同的算法简称视为同一次调用。

随后调用前节说明的ILS内核，将连续未刷新最好收益的阈值设置为$15$。若计数为$q\in\mathbb{Z}_{\geq0}$，代码在$q>15$时停止，严格改善历史最好收益后将$q$重置为$0$。因此ILS15不是固定执行$15$次扰动；在没有其他停止条件且连续未改进时，要到计数达到$16$才触发这一停止检查。适配器显式关闭ILS的时间上限，实际搜索时间还取决于纪录刷新次数及每轮交换和补全的工作量。

进入ILS还需满足适配器的保护条件：启用ILS、当前匹配收益为正，并且公开接口返回的未选受阻候选容器非空。该容器用于后续扰动候选访问，其元素数量为无量纲整数。若归约及构造后已经没有可供扰动的候选，适配器直接跳过ILS，避免对空容器执行非法随机抽样。这一保护不增加新的搜索操作，也不把没有扰动候选解释为取得了全局最优解。

ILS内部的交换、扰动、接受概率及连续未改进计数按第\ref{sec:03_03}节所述公开内核执行。完整路径中的精确标记和折叠历史必须随匹配状态保留，不能将残余图作为一个与恢复过程无关的新输入。精确标记是实现中的保护状态，并非求解器颁发的最优性证书；最终质量仍取决于实际归约、残余搜索和恢复是否正确完成。

\subsection{逆序展开与原始输入验收}

折叠可能嵌套发生：一个临时任务可以参与后续折叠，也可以被后续规则固定。因此必须按历史栈的逆序恢复，先撤销最后发生的结构变换，再恢复此前变换。每次WEF按式\eqref{s0304:wef_unfold}处理；在数学WT变换中，已选孪生代表应恢复为两个原任务。实际展开只处理本次运行真正写入历史栈的记录，不能因驱动器调用过某项规则就假定存在相应记录。固定的临时任务也沿历史继续展开，最终输出必须只含$E_0$中的原始任务编号。

适配器在展开后遍历解中的编号，检查是否超出原始任务编号范围或输出容量，并用原始输入权重重新求和。归约后贪心阶段和ILS阶段记录的内部收益仅是归约状态下的数值，不能直接当作原问题收益。尤其存在折叠时，忽略偏置或把临时任务权重当作原权重，都会破坏式\eqref{s0304:solution_offset}所表达的对应关系。

编号恢复完成仍不等于原始问题已验收。设输出编号列表为$L$，其去重后的集合为$M_0$。独立核验需回到原始JSON输入，检查
\begin{align}
 &M_0\subseteq E_0,\qquad |L|=|M_0|,
 \label{s0304:verify_ids}\\
 &\ell_0(v)=\sum_{e\in M_0}\mathbb{I}_{\{v\in R_e\}}\leq1,
 &&\forall v\in V_0,
 \label{s0304:verify_capacity}\\
 &z_{\mathrm{check}}=\sum_{e\in M_0}w_e.
 \label{s0304:verify_objective}
\end{align}
式\eqref{s0304:verify_ids}同时排除非法编号和重复输出；式\eqref{s0304:verify_capacity}恢复检查全部原始资源，而不是只检查归约后保留的资源；式\eqref{s0304:verify_objective}给出可报告的原始目标值。$\ell_0(v)$是无量纲占用计数，$z_{\mathrm{check}}$是无量纲收益。有效性要求这些检查通过，不能以程序正常返回、归约标记为精确或内部收益为正替代。原始可行性和目标求和核验也不证明归约保留了全局最优值。

适配器分别记录建图、归约、贪心、ILS以及展开和输出检查所在区间的耗时。令单调时钟读数$t_0,\ldots,t_5$依次位于这些区间的边界，则
\begin{equation}
 T_j=t_j-t_{j-1}\geq0,
 \qquad j\in\{1,2,3,4,5\}.
 \label{s0304:stage_times}
\end{equation}
各$T_j$单位均为$\mathrm{s}$，依次对应上述五个阶段。最后一个区间包含展开、编号复制及适配器中的原权重求和；在适配器之外读取原始JSON并独立验收的工作不能自动归入这一内部计时。跳过ILS时仍存在相应时钟区间，其中程序开销不代表实施了搜索。

\subsection{已有运行及其证据边界}

完整路径在$7$个已有原始实例条件上分别安排$10$个固定随机种子，共计划$70$次运行。独立核验记录确认其中$65$次有效、$5$次失败。这$7$个条件已经包含在已有实验条件中，不能作为新增独立数据集再次累计。失败运行分别出现在\texttt{ibm02}的$1$个种子和\texttt{us04}的$4$个种子；具体失败表现与适用边界在后续结果分析中讨论。

对某个条件$c$，令$\mathcal{S}_c$为计划种子集合，$|\mathcal{S}_c|=10$，令$\mathcal{V}_c\subseteq\mathcal{S}_c$为通过原始输入验收的种子集合。对$s\in\mathcal{V}_c$，记核验收益为$z_{c,s}$。只在有效运行上计算的均值是
\begin{equation}
 \overline z_c^{\mathrm{valid}}
 =\frac{1}{|\mathcal{V}_c|}
 \sum_{s\in\mathcal{V}_c}z_{c,s},
 \qquad |\mathcal{V}_c|>0.
 \label{s0304:valid_mean}
\end{equation}
若$|\mathcal{V}_c|<10$，该值必须标为有效子集均值，不能称为全部种子均值。失败对应的收益保持缺失，既不擅自记为零，也不从计划运行数中删除。只有$\mathcal{V}_c=\mathcal{S}_c$时，式\eqref{s0304:valid_mean}才同时是该条件的全部种子均值。因此，对\texttt{ibm02}和\texttt{us04}不能报告由现有有效输出直接得到的全部种子平均收益，也不能用$65$次有效运行替代$70$次计划运行来描述完成情况。

这$65$次有效结果证明的是相应输出在原始输入上通过了编号、容量和收益核验，不是所有归约规则均保持最优值的实验性证明。尤其WT的实际检测路径存在上述静态缺陷，完整驱动器被调用并不等于其每项理论规则都已在运行中得到验证。安全数学条件、实现覆盖范围和历史输出可行性必须分别陈述。

完整流程与单独调用Greedy及ILS的配置不同，历史批次之间的停止条件、热启动和CPU并发也不同。因此，已有运行可以支持逐次核验后的收益与阶段耗时陈述，却不能据此建立统一计算预算下的速度排名。对公开算法的评价需要同时保留成功恢复的结果、未完成或未通过验收的运行，以及尚未由性能实验分离的实现影响。
